\begin{proposition}
\label{prop:safe-gap}
For each $a>0$ and each $b>a$, a full orbit snapshot
at a specified $c\in(a,b)$ gives a quantitative
inverse on the constant-initial-position geodesic-image
class. Set $M=1/(4a^2)$ and choose an integer
$L>\max\{\pi/(b-a),\,2M+2\}$.
Take $c=a+\pi/L$.
Using just the full orbit snapshots at $a$ and $c$,
recovery of the initial velocity in energy normalization
has sharp loss one. In particular, for each real $s$,
\[
 \|U_t(a)-\widehat U_t(a)\|_{H^{s-1}}
 \le C(a,c)\|
       \log_{U_0}U(c)-\log_{U_0}\widehat U(c)
      \|_{H^{s+1}} .
\]
Both solutions start at the same constant target point $U_0$.
The velocity difference is taken in $T_{U_0}\mathbb H^2$.
\end{proposition}

\begin{proof}
The Riemannian logarithm turns a geodesic through $U_0$
into a straight line in $T_{U_0}\mathbb H^2$;
therefore $Z=\log_{U_0}U$ satisfies the scalar
equation componentwise, also when comparing two
different geodesics. On frequency $k\ge1$, rescale
by $\sqrt t$ and normalize the initial slope to $k$.
Writing the resulting solution as
$z=\rho\sin\varphi$, $z'=k\rho\cos\varphi$ gives
\[
 \varphi'=k+\frac{W(t)}{k}\sin^2\varphi,\qquad
 \frac{\rho'}{\rho}=-\frac{W(t)}{2k}\sin(2\varphi),
 \qquad W(t)=\frac1{4t^2}.
\]
The starting values are $\rho(a)=1$, $\varphi(a)=0$.
With $d=\pi/L$, the correction to the free phase at $c$
is strictly between zero and $Md/k$. If $k>M$,
this is less than the spacing $d$ of the free phases
modulo $\pi$. No mode can reach a zero. If $k\le M$,
the phase is between zero and $(k+M/k)d\le 2Md<\pi$.
There are thus no finite-frequency zeros.

For large $k$ the same phase bounds hold at every time
in the window, so
\[
 \varphi(c)-kd
 =\frac1k\int_a^c W(t)\sin^2(k(t-a))\,dt+O(k^{-2})
 =\frac1{2k}\int_a^c W(t)\,dt+O(k^{-2}).
\]
The second equality uses integration by parts.
The amplitude is $1+O(k^{-1})$.
On multiples of $L$, the multiplier is therefore of
order $k^{-1}$ with a nonzero leading coefficient.
On large nonmultiples it is bounded away from zero.
The zero-mode transfer is $a\log(c/a)$
for the unscaled velocity. This proves the estimate.
Taking single high modes on one target geodesic shows
sharpness for the pair $a,c$.
\end{proof}
